\section{Analysis on Multi-Agent Systems with Open-Source Models}
\label{sec:analysis_on_opensource_models}

In this section, we also provide the analysis of failure modes on MetaGPT and ChatDev frameworks where the underlying LLMs are open-source models. In particular, we chose to use \verb|Qwen2.5-Coder-32B-Instruct| \cite{hui2024qwen2} and \verb|CodeLlama-7b-Instruct-hf| \cite{roziere2023code} models for these two frameworks. The analysis of failure modes is shown in Table \ref{tab:open_source_models}. This new analysis reveals two key findings:
\begin{itemize}
    \item There is a significant performance difference between the two open-source models. \verb|Qwen2.5-Coder-32B-Instruct| is substantially more robust than \verb|CodeLlama-7b-Instruct-hf| on these tasks, exhibiting far fewer failures overall.
    \item  Both open-source models show a higher frequency of failures compared to the leading closed-source models analyzed in our paper (GPT-4o and Claude-3). This suggests a performance gap and highlights important areas for future improvement in open-source models for multi-agent tasks.
\end{itemize}

\begin{table*}[ht!]
    \small
    \renewcommand{\arraystretch}{1.3}
    \setlength{\tabcolsep}{4pt}
    \centering
    \caption{Failure Mode Occurrences in 400 Traces with Open-Source Models. Results are grouped by model family (Qwen vs.\ CodeLlama) and development framework (ChatDev vs.\ MetaGPT).}
    \label{tab:failure-modes-open-source}
    \begin{tabular}{p{2.5cm} c c c c}
        \toprule
        \textbf{Failure Mode} & \multicolumn{2}{c}{\textbf{Qwen}} & \multicolumn{2}{c}{\textbf{CodeLlama}} \\
        \cmidrule(lr){2-3} \cmidrule(lr){4-5}
         & \textbf{ChatDev} & \textbf{MetaGPT} & \textbf{ChatDev} & \textbf{MetaGPT} \\
        \midrule
        1.1 & 35 & 12 & 76 & 94 \\
        \hline
        1.2 & 4 & 1 & 45 & 12 \\
        \hline
        1.3 & 96 & 35 & 97 & 99 \\
        \hline
        1.4 & 1 & 0 & 46 & 23 \\
        \hline
        1.5 & 94 & 3 & 97 & 76 \\
        \hline
        2.1 & 2 & 0 & 50 & 9 \\
        \hline
        2.2 & 1 & 4 & 16 & 15 \\
        \hline
        2.3 & 9 & 0 & 76 & 57 \\
        \hline
        2.4 & 0 & 0 & 2 & 0 \\
        \hline
        2.5 & 2 & 12 & 42 & 40 \\
        \hline
        2.6 & 20 & 16 & 93 & 18 \\
        \hline
        3.1 & 1 & 47 & 25 & 26 \\
        \hline
        3.2 & 16 & 51 & 67 & 55 \\
        \hline
        3.3 & 12 & 32 & 69 & 56 \\
        \hline
    \end{tabular}
\label{tab:open_source_models}
\end{table*}
